\documentclass[border=12pt,tikz]{standalone}
\usepackage{amsmath,amssymb}
\usepackage{CJKutf8}
\usetikzlibrary{arrows.meta,calc,decorations.markings}

\begin{document}
\begin{CJK}{UTF8}{ipxm}
\begin{tikzpicture}[
    >=Stealth,
    font=\small,
    every node/.style={inner sep=2pt},
]

\def\Ra{2.0}  % R1 (small)
\def\Rb{3.2}  % R2 (large)

% Angle of ray OP1 from vertical
\pgfmathsetmacro{\thDeg}{atan(\Ra/\Rb)}  % arctan(R1/R2)
\pgfmathsetmacro{\raylen}{sqrt(\Ra*\Ra+\Rb*\Rb)}

% ============================================================
%  TWO CONCENTRIC CIRCLES
% ============================================================
\draw[red!50!black, thick] (0,0) circle (\Ra);
\node[red!50!black, font=\footnotesize] at ({-\Ra*cos(45)},{-\Ra*sin(45)-0.25}) {$S(R_1)$};

\draw[blue!50!black, thick] (0,0) circle (\Rb);
\node[blue!50!black, font=\footnotesize] at ({-\Rb*cos(45)-0.1},{-\Rb*sin(45)-0.25}) {$S(R_2)$};

% Origin O
\filldraw (0,0) circle (2pt);
\node[below left=3pt] at (0,0) {$O$};

% ============================================================
%  TANGENT LINES
% ============================================================

% Blue: horizontal tangent at y = R2
\draw[blue, very thick] (-1.0,\Rb) -- (4.0,\Rb);
\node[blue, above=2pt, font=\footnotesize] at (-0.5,\Rb) {$\Pi_{R_2}$};

% Red: vertical tangent at x = R1
\draw[red!70!black, very thick] (\Ra,-0.8) -- (\Ra,4.0);
\node[red!70!black, right=3pt, font=\footnotesize] at (\Ra,3.8) {$\Pi_{R_1}$};

% ============================================================
%  INTERSECTION POINT P1
% ============================================================
\coordinate (P1) at (\Ra,\Rb);
\filldraw[black] (P1) circle (3pt);
\node[above right=3pt, font=\normalsize] at (P1) {$P_1$};
\node[right=15pt, font=\footnotesize, gray] at (P1) {$(R_1, R_2)$};

% ============================================================
%  PROJECTION RAY from O through P1
% ============================================================
\draw[gray, dashed] (0,0) -- (P1);

% Points on each sphere (same ray, same angle theta)
\pgfmathsetmacro{\QBx}{\Rb*\Ra/\raylen}  % on S(R2)
\pgfmathsetmacro{\QBy}{\Rb*\Rb/\raylen}
\pgfmathsetmacro{\QAx}{\Ra*\Ra/\raylen}  % on S(R1)
\pgfmathsetmacro{\QAy}{\Ra*\Rb/\raylen}

\filldraw[blue!70!black] (\QBx,\QBy) circle (3pt);
\filldraw[red!60!black] (\QAx,\QAy) circle (3pt);

% ============================================================
%  THICK ARCS from north poles to projection points
% ============================================================

% Blue arc: from (0, R2) to Phi_{R2}(P1) along S(R2)
% Angle from vertical = thDeg, arc from 90 deg to (90-thDeg) deg
\draw[blue, line width=3pt, opacity=0.7,
  -{Stealth[length=8pt,width=6pt]}]
  (90:\Rb) arc (90:{90-\thDeg}:\Rb);
\node[blue, font=\normalsize, left=8pt] at ({0.5*\Rb*sin(\thDeg/2)},{\Rb*cos(\thDeg/2)})
  {$R_2\theta$};

% Red arc: from (0, R1) to Phi_{R1}(P1) along S(R1)
% Same angle thDeg, but on smaller circle
% North pole of S(R1) at (0, R1)
\filldraw[red!60!black] (0,\Ra) circle (2pt);

\draw[red!70!black, line width=3pt, opacity=0.7,
  -{Stealth[length=8pt,width=6pt]}]
  (90:\Ra) arc (90:{90-\thDeg}:\Ra);
\node[red!70!black, font=\normalsize, left=6pt] at ({0.3*\Ra*sin(\thDeg/2)},{\Ra*cos(\thDeg/2)})
  {$R_1\theta$};

% ============================================================
%  ANGLE theta
% ============================================================
\draw[gray, thin] (0,0.8) arc (90:{90-\thDeg}:0.8);
\node[gray, font=\footnotesize, above=1pt] at ({0.8*sin(\thDeg/2)},{0.8*cos(\thDeg/2)})
  {$\theta$};

% ============================================================
%  Contact points
% ============================================================
\filldraw[blue] (0,\Rb) circle (2.5pt);
\node[blue, left=4pt, font=\footnotesize] at (0,\Rb) {$P_{R_2}$};

\filldraw[red!70!black] (\Ra,0) circle (2.5pt);
\node[red!70!black, below=4pt, font=\footnotesize] at (\Ra,0) {$P_{R_1}$};

% R1, R2 braces
\draw[red!70!black, thick, decorate, decoration={brace, amplitude=4pt, mirror, raise=3pt}]
  (0,-0.1) -- (\Ra,-0.1);
\node[red!70!black, below=8pt, font=\footnotesize] at (\Ra/2,0) {$R_1$};

\draw[blue, thick, decorate, decoration={brace, amplitude=4pt, raise=3pt}]
  (-0.1,0) -- (-0.1,\Rb);
\node[blue, left=8pt, font=\footnotesize] at (0,\Rb/2) {$R_2$};

% ============================================================
%  KEY MESSAGE
% ============================================================
\node[font=\normalsize, align=left, text width=9cm] at (2.5,-2.5)
  {同じ背景座標の点 $P_1$ への射影線は同じ角度 $\theta$ を成す．\\
   弧の長さ（主観空間での距離）は $R_1\theta$ と $R_2\theta$．\\
   \textbf{比は $R_1 : R_2$}（曲率半径の比に等しい）．};

\end{tikzpicture}
\end{CJK}
\end{document}
